                        IMC 2019, Blagoevgrad, Bulgaria


                                         Day 1, July 30, 2019




Problem 1. Evaluate the product
                                                ∞
                                                Y (n3 + 3n)2
                                                                    .
                                                n=3
                                                          n6 − 64
         Proposed by Orif Ibrogimov, ETH Zurich and National University of Uzbekistan and
        Karen Keryan, Yerevan State University and American University of Armenia, Yerevan

Hint: Telescoping product.


Solution. Let
                                                       (n3 + 3n)2
                                                an =              .
                                                        n6 − 64
     Notice that
                        (n3 + 3n)2                     n2 (n2 + 3)2
                an =                  =
                     (n3 − 8)(n3 + 8)   (n − 2)(n2 + 2n + 4) · (n + 2)(n2 − 2n + 4)
                       n       n       n2 + 3       n2 + 3
                   =       ·       ·            ·             .
                     n − 2 n + 2 (n − 1)2 + 3 (n + 1)2 + 3
Hence, for N ≥ 3 we have
          N             N
                                     !   N
                                                      !     N
                                                                              !   N
                                                                                                        !
                                                                        2                 2
          Y             Y    n           Y    n             Y     n +3            Y      n +3
                an =
          n=3           n=3
                            n−2          n=3
                                             n+2           n=3
                                                               (n − 1)2 + 3       n=3
                                                                                      (n + 1)2 + 3
                     N (N − 1)        3·4         N2 + 3   32 + 3
                   =           ·                · 2      ·
                        1·2      (N + 1)(N + 2) 2 + 3 (N + 1)2 + 3
                     72        N (N − 1)(N 2 + 3)
                   =    ·                            
                      7 (N + 1)(N + 2) (N + 1)2 + 3
                     72          (1 − N1 )(1 + N32 )
                   =    ·                                   ,
                      7 (1 + N1 )(1 + N2 ) (1 + N1 )2 + N32
so
           ∞                  N
                                                                                              !
           Y                  Y                 72         (1 − N1 )(1 + N32 )                        72
                an = lim            an = lim       ·                                              =      .
                                                7 (1 + N1 )(1 + N2 ) (1 + N1 )2 + N32
                                                                                      
          n=3
                       N →∞
                              n=3
                                         N →∞                                                         7

Problem 2. A four-digit number            Y EAR is called very good if the system

                                          Y x + Ey + Az + Rw = Y
                                          Rx + Y y + Ez + Aw = E
                                          Ax + Ry + Y z + Ew = A
                                          Ex + Ay + Rz + Y w = R

of linear equations in the variables x, y, z and w has at least two solutions. Find all very good
YEARs in the 21st century.
    (The 21st century starts in 2001 and ends in 2100.)
                                           Proposed by Tomá² Bárta, Charles University, Prague

                                                           1
                                                                          
                                                             Y   E   A   R
                                                           R    Y   E   A
Hint: If the solution of the system is not unique then det                 = 0.
                                                           A    R   Y   E
                                                             E   A   R   Y

Solution. Let us apply row transformations to the augmented matrix of the system to nd its
rank. First we add the second, third and fourth row to the rst one and divide by Y +E +A+R
to get                                                                     
                  1 1 1 1 1               1     1        1        1       1
                R Y E A E  0 Y − R E − R A − R E − R
                A R Y E A ∼ 0 R − A Y − A E − A
                                                                           
                                                                          0 
                  E A R Y R               0 A−E R−E Y −E R−E
                                                                           
                      1         1             1            1            1
                    0       Y −R          E−R          A−R           E − R
                 ∼                                                         
                    0       R−A           Y −A         E−A             0 
                      0 A−E+Y −R              0    Y −E+A−R             0
Let us rst omit the last column and look at the remaining 4 × 4 matrix. If E 6= R, the rst
and second rows are linearly independent, so the rank of the matrix is at least 2 and rank of
the augmented 4 × 5 matrix cannot be bigger than rank of the 4 × 4 matrix due to the zeros in
the last column.
   IF E = R, then we have three zeros in the last column, so rank of the 4×5 matrix cannot be
bigger than rank of the 4 × 4 matrix. So, the original system has always at least one solution.
   It follows that the system has more than one solution if and only if the 4 × 4 matrix (with
the last column omitted) is singular. Let us rst assume that E 6= R. We apply one more
transform to get
                                                                                          
      1                    1                      1                      1
    0                  Y −R                   E−R                    A−R                  
  ∼0 (R − A)(E − R) − (Y − R)(Y − A)
                                                                                           
                                                  0     (E − A)(E − R) − (A − R)(Y − A)
      0             A−E+Y −R                      0              Y −E+A−R

Obviously, this matrix is singular if and only if A − E + Y − R = 0 or the two expressions in
the third row are equal, i.e.

RE − R2 − AE + AR − Y 2 + RY + AY − AR = E 2 − AE − ER + AR − AY + RY + A2 − AR

                                  0 = (E − R)2 + (A − Y )2 ,
but this is impossible if E 6= R. If E = R, we have
                                                                    
             1       1            1         1           1  1   1   1
           0      Y −R           0       A−R     ∼ 0 Y − R
                                                              0 A − R
        ∼                                                             .
           0      R−A         Y −A       R − A  0 R − A Y − A R − A 
             0 A + Y − 2R         0    Y + A − 2R       0 A−R  0 Y −R

If A = Y , this matrix is singular. If A 6= Y , the matrix is regular if and only if (Y − R)2 6=
(A − R)2 and since Y 6= A, it means that Y − R 6= −(A − R), i.e. Y + A 6= 2R. We conclude
that YEAR is very good if and only if
     1. E 6= R and A + Y = E + R, or
     2. E = R and Y = A, or
     3. E = R, A 6= Y and Y + A = 2R.


                                               2
   We can see that if Y = 2, E = 0, then the very good years satisfying 1 are A+2 = R 6= 0, i.e.
2002, 2013, 2024, 2035, 2046, 2057, 2068, 2079, condition 2 is satised for 2020 and condition 3
never satised.

Problem 3. Let     f : (−1, 1) → R be a twice dierentiable function such that

                               2f 0 (x) + xf 00 (x) ≥ 1 for x ∈ (−1, 1).

Prove that                                  Z 1
                                                           1
                                                xf (x) dx ≥ .
                                             −1            3
          Proposed by Orif Ibrogimov, ETH Zurich and National University of Uzbekistan and
           Karim Rakhimov, Scuola Normale Superiore and National University of Uzbekistan

Hint:   2f 0 (x) + xf 00 (x) is the second derivative of a certain function.

Solution. Let
                                                                 x2
                                          g(x) = xf (x) −           .
                                                                 2
Notice that
                                  g 00 (x) = 2f 0 (x) + xf 00 (x) − 1 ≥ 0,
so g is convex. Estimate g by its tangent at 0: let g 0 (0) = a, then

                                      g(x) = g(0) + g 0 (0)x = ax

and therefore
                  Z 1                Z 1                         Z 1
                                             x2                               x2      1
                       xf (x) dx =    g(x) +      dx ≥                    ax +      dx = .
                    −1             −1        2                         −1      2        3

Problem 4. Dene the sequence         a0 , a1 , . . . of numbers by the following recurrence:

                a0 = 1,    a1 = 2,    (n + 3)an+2 = (6n + 9)an+1 − nan             for n ≥ 0.

Prove that all terms of this sequence are integers.
                                          Proposed by Khakimboy Egamberganov, ICTP, Italy

Hint: Determine the generating function       an x n .
                                           P

Solution. Take the generating function of this sequence

                                                      ∞
                                                      X
                                            f (x) =         an x n .
                                                      n=0

It is easy to see that the sequence is increasing and

              an+1   (6n + 3)an − (n − 1)an−1   6n + 3                                  an+1
                   =                          <                          ⇒    lim sup        ≤ 6.
               an           (n + 2)an            n+2                           n→∞       an
So the generating function converges in some neighbourhood of 0. Then, we have
               ∞                ∞                      ∞                    ∞
             X
                      n
                              X
                                          n+2
                                                     X    6n + 9       n+2
                                                                           X     n
f (x) = 1+2x+     an x = 1+2x+     an+2 x     = 1+2x+            an+1 x −            an xn+2 .
              n=2              n=0                    n=0
                                                          n + 3            n=0
                                                                               n + 3

                                                      3
                       ∞                                                   ∞
                       X 6n + 9                                            X        n
Let f1 (x) =                           an+1 x   n+2
                                                       and f2 (x) =                     an xn+2 . Then
                       n=0
                              n+3                                             n=0
                                                                                  n + 3
                   ∞
                   X                                              ∞
                                                                  X                                   ∞
                                                                                                      X
           0
(xf1 (x)) =                  (6n+9)an+1 x      n+2
                                                      = 6x    2
                                                                        (n+1)an+1 x +3x    n
                                                                                                            an+1 xn+1 = 6x2 f 0 (x)+3x(f (x)−1)
                       n=0                                        n=0                                 n=0
and
                       ∞
                       X                             ∞
                                                     X                                  ∞
                                                                                        X
               0
(xf2 (x)) =                   nan x   n+2
                                            =x   2
                                                           (n + 1)an x − xn         2
                                                                                               an xn = x2 (xf (x))0 − x2 f (x) = x3 f 0 (x).
                       n=0                           n=0                                 n=0
Using this relations, we arrive at the following dierential equation for f :
           (xf (x))0 = 1 + 4x + (xf1 (x))0 − (xf2 (x))0 = 1 + x + (6x2 − x3 )f 0 (x) + 3xf (x)
or, equivalently,
                                       (x3 − 6x2 + x)f 0 (x) + (1 − 3x)f (x) − 1 − x = 0.
So, we need solve this dierential equation in some suciently smaller neighbourhood of 0. We
know that f (0) = 1 and we need a neighbourhood of 0 such that x2 − 6x + 1 > 0. Then
                                          1 − 3x                   1+x
                          f 0 (x) +      2
                                                    f (x) =      2
                                    x(x − 6x + 1)             x(x − 6x + 1)
                                                           x
for x 6= 0. So the integral multiplier is µ(x) = √                  and
                                                       2
                                                      x − 6x + 1
                                                         x + x2
                                    (f (x)µ(x))0 =                 3 ,
                                                   (x2 − 6x + 1) 2
so                                            √ 2                        √
                            1−x              1   x − 6x + 1        1 − x − x2 − 6x + 1
             f (x) =    √                 −                     =                        .
                       2 x2 − 6x + 1 2                x                      2x
We found the generating function of (an ) in some neighbourhood
                                                           √             of 0, which x2 − 6x + 1 > 0.
                                                 1 − x − x2 − 6x + 1                    √
So our series uniformly converges to f (x) =                              in |x| < 3 − 2 2.
                                                             2x
    Instead of computing the coecients of the Taylor series of f (x) directly, we will nd
another recurrence relation for (an ). It is easy to see that f (x) satises the quadratic equation
xt2 − (1 − x)t + 1 = 0. So
                                    xf (x)2 − (1 − x)f (x) + 1 = 0.
Then
    ∞
                       !2          ∞                 ∞                              ∞ n
                                                                                                                !                 ∞
    X                              X                 X                              X X                                           X
                   n                           n                   n+1                                                  n+1
x         an x              +1 =            an x −          an x          ⇒                           ak an−k       x         =         (an+1 −an )xn+1
    n=0                            n=0               n=0                             n=0        k=0                               n=0

and from here, we get
                                                                                   n
                                                                                   X
                                                            an+1 = an +                  ak an−k .
                                                                                   k=0
If a0 , a1 , ..., an be integers, then an+1 is also integer. We know that a0 = 1, a1 = 2 are integer
numbers, so all terms of the sequence (an ) are integers by induction.
Problem 5.   Determine whether there exist an odd positive integer n and n × n matrices A
and B with integer entries, that satisfy the following conditions:
     (1) det(B) = 1;
     (2) AB = BA;
     (3) A4 + 4A2 B 2 + 16B 4 = 2019I .
   (Here I denotes the n × n identity matrix.)
            Proposed by Orif Ibrogimov, ETH Zurich and National University of Uzbekistan

                                                                               4
Hint: Consider the determinants modulo      4.

     Remark. The proposed solution was more complicated and involved; during the contest it
turned out that a signcantly simplied solution exists  which we now provide below.

     Solution 1. We show that there are no such matrices.
     Notice that A4 + 4A2 B 2 + 16B 4 can factorized as

                  A4 + 4A2 B 2 + 16B 4 = (A2 + 2AB + 4B 2 )(A2 − 2AB + 4B 2 ).

Let C = A2 + 2AB + 4B 2 and D = A2 − 2AB + 4B 2 be the two factors above. Then

          det C · det D = det(CD) = det(A4 + 4A2 B 2 + 16B 4 ) = det(2019I) = 2019n .

   The matrices C, D have integer entries, so their determinants are integers. Moreover, from
C ≡ D (mod 4) we can see that

                                    det C ≡ det D    (mod 4).

This implies that det C·det D ≡ (det C)2 (mod 4), but this is a contradiction because 2019n ≡ 3
(mod 4) is a quadratic nonresidue modulo 4.

     Solution 2. Notice that


                           A4 ≡ A4 + 4A2 B 2 + 16B 4 = 2019I    mod 4

so
                       (det A)4 = det A4 ≡ det(2109I) = 2019n    (mod 4).
But 2019n ≡ 3 is a quadratic nonresidue modulo 4, contradiction.




                                                 5
